\section*{Chapitre \ref{ch:qm:covariance-estimation-and-random-matrices} --- Estimation de la covariance et matrices aléatoires}

\begin{solution}{exo:qm:covariance-estimation-and-random-matrices:1}
$q = 0.4$~: $\lambda_\pm = (1 \pm \sqrt{0.4})^2 = 0.135$ et $2.665$ (pour une matrice de corrélation, $\sigma^2 = 1$).
\end{solution}

\begin{solution}{exo:qm:covariance-estimation-and-random-matrices:2}
$q = 0.5$, soit une volatilité vraie d'environ $5\%/(1 - 0.5) = 10\%$.
\end{solution}

\begin{solution}{exo:qm:covariance-estimation-and-random-matrices:3}
Avec $w = \Sigma^{-1}\mathbf 1/c$, $c = \mathbf 1^\top\Sigma^{-1}\mathbf 1$~: $w^\top\Sigma w = \mathbf 1^\top\Sigma^{-1}\Sigma\Sigma^{-1}\mathbf 1/c^2 = c/c^2 = 1/c$.
\end{solution}

\begin{solution}{exo:qm:covariance-estimation-and-random-matrices:4}
Des pics au-dessus de $1 + \sqrt{0.25} = 1.5$. Un pic de 3 apparaît en $3(1 + 0.25/2) = 3.375$, au-dessus de l'edge $(1.5)^2 = 2.25$.
\end{solution}

\begin{solution}{exo:qm:covariance-estimation-and-random-matrices:5}
$\delta mI + (1 - \delta)U\Lambda U^\top = U(\delta mI + (1 - \delta)\Lambda)U^\top$. Chaque valeur propre est attirée vers $m$ dans la même proportion, de sorte que les petites valeurs propres, qui sont trop petites,
et les grandes, qui sont trop grandes, sont corrigées par un seul nombre~; la correction oracle n'est pas linéaire en $\lambda$, alors que le \omterm{def:qm:covariance-estimation-and-random-matrices:nonlinear}{rétrécissement non linéaire} la choisit valeur propre par valeur propre.
\end{solution}

\begin{solution}{exo:qm:covariance-estimation-and-random-matrices:6}
$1/(1 - q) \le 1.1$ exige $q \le 0.091$, d'où $T \ge 200/0.091 = 2\,200$ jours, soit environ 8.7 ans~; sur une telle durée, la covariance elle-même change.
\end{solution}

\begin{solution}{exo:qm:covariance-estimation-and-random-matrices:7}
\texttt{bias\_vs\_q()}~: 1.11, 1.26, 1.43, 1.67, 2.05, 2.51, 3.38, 5.21 et 10.05 pour $q = 0.1$ à 0.9, contre 1.11, 1.25, 1.43, 1.67, 2.00, 2.50, 3.33, 5.00 et 10.09.
\end{solution}

\begin{solution}{exo:qm:covariance-estimation-and-random-matrices:8}
La volatilité en échantillon d'un \omterm{def:qm:covariance-estimation-and-random-matrices:pca}{portefeuille de variance minimale} est la plus faible pour l'\omterm{def:qm:estimation:estimator}{estimateur} dont les petites valeurs propres sont les plus sous-estimées, c'est-à-dire le plus bruité~: dans le chapitre, la
\omterm{def:qm:covariance-estimation-and-random-matrices:cov}{matrice de covariance empirique} brute affiche la plus faible valeur en échantillon (5.7~\%) et le pire risque réalisé (9.6~\%). Il faut sélectionner les \omterm{def:qm:estimation:estimator}{estimateurs} selon leur risque hors échantillon, ou selon le ratio entre risque réalisé et risque prédit.
\end{solution}

\begin{solution}{pb:qm:covariance-estimation-and-random-matrices:1}
\textbf{1.} Prédit 5.71~\%, réalise 9.58~\% contre la valeur vraie (9.61~\% sur l'année suivante)~; le véritable optimum est 7.44~\%.
\textbf{2.} $1/(1 - 0.4) = 1.67$~; la simulation donne 1.68.
\textbf{3.} 0.135 et 2.665~; sur du bruit pur, les valeurs propres empiriques vont de 0.144 à 2.565.
\textbf{4.} Le marché (51.4 contre 52.9 dans la population) et quatre secteurs (2.96 à 3.68 contre 2.83 à 2.99)~: les valeurs propres sectorielles sont poussées vers le haut par le bruit, comme
le prédit $\ell(1 + q/(\ell - 1))$.
\textbf{5.} Un pic égal à $1 + \sqrt{0.4} = 1.63$ fois la variance du bruit.
\textbf{6.} 0.034 vers l'identité, pour un risque réalisé de 9.14~\%~; 0.178 vers la corrélation constante, également 9.14~\% mais avec une prédiction honnête (7.64~\%).
\textbf{7.} \omterm{def:qm:covariance-estimation-and-random-matrices:nonlinear}{Rétrécissement non linéaire} 8.25~\% (prédit 7.36~\%)~; écrêtage 8.05~\% (prédit 7.27~\%).
\textbf{8.} 7.74~\% (prédit 6.65~\%)~: la vérité est un modèle à six facteurs, que six composantes principales restituent, et le reste diagonal est exact.
\textbf{9.} 1.68, 1.51, 1.19, 1.12, 1.11 et 1.17.
\textbf{10.} L'écrêtage et le \omterm{def:qm:covariance-estimation-and-random-matrices:nonlinear}{rétrécissement non linéaire} ont les prédictions les plus honnêtes (ratios 1.11 et 1.12)~; le \omterm{def:qm:covariance-estimation-and-random-matrices:factor}{modèle factoriel} construit le meilleur portefeuille (7.74~\%).
\textbf{11.} Sa volatilité vraie vaut dix fois sa prédiction (10.05 simulée, 10.09 par la formule).
\textbf{12.} $q \le 0.091$~: 2\,200 jours.
\textbf{13.} Il pondère chaque direction par l'inverse de sa variance estimée, et les variances les plus sous-estimées sont du bruit~; l'optimiseur se concentre précisément là où l'estimation est
la moins fiable.
\textbf{14.} Elle interdit les grandes positions compensatoires qui exploitent le bruit, ce qui agit comme une forme de rétrécissement (au prix d'exclure les ventes à découvert réellement utiles).
\textbf{15.} Les actions partagent un facteur de marché~; la cible identité suppose qu'elles ne sont pas corrélées, de sorte qu'un fort rétrécissement vers elle serait très erroné et que l'intensité optimale reste minime (0.034).
\textbf{16.} \omterm{def:qm:covariance-estimation-and-random-matrices:nonlinear}{Rétrécissement non linéaire} ou \omterm{def:qm:covariance-estimation-and-random-matrices:factor}{modèle factoriel} dont le reste est rétréci, contrôlé par écrêtage~; pas la covariance empirique brute.
\textbf{17.} Suivre la volatilité réalisée comparée à la volatilité prédite des portefeuilles en production sur une fenêtre glissante, ainsi que la stabilité des valeurs propres au-dessus de l'edge.
\textbf{18.} Des queues plus lourdes élargissent la dispersion des valeurs propres et favorisent les \omterm{def:qm:estimation:estimator}{estimateurs} robustes ou fondés sur les rangs~; une covariance variable raccourcit l'historique utilisable, ce qui augmente $q$ et rend le nettoyage
plus important.
\textbf{19.} Résultat nommé~: \emph{le portefeuille qui promettait 6~\%}~: pour $q = 0.4$, le \omterm{def:qm:covariance-estimation-and-random-matrices:pca}{portefeuille de variance minimale} empirique réalise 1.68 fois la volatilité qu'il prédit (5.71~\% annoncés, 9.58~\%
réalisés), contre $1/(1 - q) = 1.67$~; le \omterm{def:qm:covariance-estimation-and-random-matrices:linear}{rétrécissement linéaire} vers l'identité donne 1.51, vers la corrélation constante 1.19, le \omterm{def:qm:covariance-estimation-and-random-matrices:nonlinear}{rétrécissement non linéaire} 1.12, l'écrêtage 1.11 et un modèle à six facteurs 1.17.
\textbf{20.} Parce qu'il cherche les directions de risque estimé le plus faible, qui sont celles où l'erreur d'estimation a fait paraître le risque le plus petit.
\end{solution}

\begin{solution}{iq:qm:covariance-estimation-and-random-matrices:1}
L'optimiseur sélectionne les directions dont la variance empirique est la plus petite, et ce sont, de façon disproportionnée, des directions où le bruit a fait baisser l'estimation~; en grande dimension, la
variance en échantillon est biaisée à la baisse d'environ $1 - q$ et la variance vraie à la hausse de $1/(1 - q)$, $q = N/T$.
\iqlookfor{Sélection sur le bruit et mise à l'échelle en $1/(1 - q)$.}
\end{solution}

\begin{solution}{iq:qm:covariance-estimation-and-random-matrices:2}
La densité limite des valeurs propres d'une \omterm{def:qm:covariance-estimation-and-random-matrices:cov}{matrice de covariance empirique} de bruit pur, portée par $[\sigma^2(1 - \sqrt q)^2, \sigma^2(1 + \sqrt q)^2]$. Sur une matrice de corrélation, les valeurs propres
situées au-dessus de l'edge supérieur portent de l'information (marché, secteurs)~; celles qui se trouvent dans le corps du spectre sont indiscernables du bruit et peuvent être nettoyées (écrêtées ou rétrécies).
\iqlookfor{Les edges, et la lecture corps du spectre contre pics.}
\end{solution}

\begin{solution}{iq:qm:covariance-estimation-and-random-matrices:3}
Une combinaison convexe de la \omterm{def:qm:covariance-estimation-and-random-matrices:cov}{matrice de covariance empirique} et d'une cible structurée, avec l'intensité qui minimise la perte de Frobenius espérée, estimée à partir des données. L'intensité croît avec le
bruit d'échantillonnage de $S$ (grand $N/T$, queues épaisses) et décroît avec la distance entre $S$ et la cible.
\iqlookfor{Compromis biais-variance et ce qui détermine l'intensité.}
\end{solution}

\begin{solution}{iq:qm:covariance-estimation-and-random-matrices:4}
$q = 1\,000/750 > 1$~: la \omterm{def:qm:covariance-estimation-and-random-matrices:cov}{matrice de covariance empirique} est singulière. Utiliser un \omterm{def:qm:covariance-estimation-and-random-matrices:factor}{modèle factoriel} (statistique ou fondamental) pour la partie commune, avec un reste diagonal ou rétréci, ou un \omterm{def:qm:covariance-estimation-and-random-matrices:nonlinear}{rétrécissement non linéaire}
conçu pour $N > T$~; vérifier à l'aide du risque réalisé des portefeuilles qu'il produit.
\iqlookfor{Reconnaître la singularité et exploiter la structure.}
\end{solution}

\begin{solution}{iq:qm:covariance-estimation-and-random-matrices:5}
Celles dont les valeurs propres dépassent l'edge supérieur de la \omterm{thm:qm:covariance-estimation-and-random-matrices:mp}{loi de Marchenko--Pastur} pour le $q$ de l'échantillon (après prise en compte de la variance qu'elles prélèvent sur le corps du spectre)~: typiquement le marché et une poignée de
secteurs~; les autres ne se distinguent pas du bruit avec autant de données.
\iqlookfor{L'edge comme critère, et le seuil des pics.}
\end{solution}

\begin{solution}{iq:qm:covariance-estimation-and-random-matrices:6}
Un \omterm{def:qm:estimation:estimator}{estimateur} qui conserve les vecteurs propres empiriques et ne modifie que les valeurs propres, de sorte qu'il se transforme de manière cohérente avec toute rotation des données. Le meilleur de cette classe utilise les valeurs propres oracle
$u_i^\top\Sigma u_i$, soit la variance vraie de chaque portefeuille propre empirique~; le \omterm{def:qm:covariance-estimation-and-random-matrices:nonlinear}{rétrécissement non linéaire} les estime à partir du spectre empirique grâce à la théorie des matrices aléatoires.
\iqlookfor{Équivariance, oracle et manière de l'approcher.}
\end{solution}
